\documentclass[11pt]{article}

\usepackage[utf8]{inputenc}
\usepackage[T1]{fontenc}
\usepackage[margin=1in]{geometry}
\usepackage{booktabs}
\usepackage{hyperref}
\usepackage{xcolor}
\usepackage{enumitem}

\hypersetup{colorlinks=true, urlcolor=blue, linkcolor=black}

\title{A Control Dataset for Behavioural Testing of Sentiment Models}

\begin{document}
\maketitle

\section{The property I want to test}

A sentiment classifier can score high on a held-out test set and still fail
on inputs a person would find trivial. The reason is that accuracy on one
distribution hides how the model behaves when you change one thing at a time.
Ribeiro et al.\ (2020) make this case directly: held-out datasets carry the
same biases as the training data, so real-world performance gets
overestimated, and a single accuracy number tells you nothing about
\emph{where} a model breaks.\footnote{Marco Tulio Ribeiro, Tongshuang Wu,
Carlos Guestrin, Sameer Singh. \emph{Beyond Accuracy: Behavioral Testing of
NLP Models with CheckList}. ACL 2020.
\url{https://arxiv.org/abs/2005.04118}}

Their answer, borrowed from software testing, is to test capabilities one by
one with inputs you control. My dataset does exactly that for three
behaviours a sentiment model should handle but frequently does not:

\begin{itemize}[noitemsep]
  \item reading negation rather than keying on a positive word,
  \item ignoring a swapped name or place that carries no sentiment,
  \item ignoring neutral filler appended to a clear sentence.
\end{itemize}

Each behaviour gets its own subset, and each subset changes exactly one
factor so that any failure points at one cause.

\section{The three subsets}

The dataset follows the three CheckList test types. A Minimum Functionality
Test (MFT) is a set of simple labelled examples that target one behaviour, in
the spirit of a unit test. An Invariance test (INV) applies a change that
should not move the prediction. A Directional Expectation test (DIR) applies
a change whose effect on the prediction is known in advance. The paper
defines all three and reports failure rates for commercial and research
models on a sentiment task, which is what makes sentiment a clean domain to
copy.

\subsection{Subset 1 -- Negation (MFT)}

I take a neutral noun and a sentiment-laden adjective and negate them with
three patterns of increasing difficulty:
\begin{center}
\texttt{The \{thing\} was not \{adjective\}.} \\
\texttt{Although they tried, the \{thing\} was not \{adjective\}.} \\
\texttt{The \{thing\} was anything but \{adjective\}.}
\end{center}
When the adjective is positive, the negated sentence is negative; when it is
negative, the negated sentence reads positive or neutral. The label is
therefore fixed by construction, which is what lets this subset carry gold
labels. A model that detects the word ``great'' and ignores the negation
gets the whole positive-adjective half wrong.

\begin{table}[h]
\centering
\begin{tabular}{ll}
\toprule
Text & Label \\
\midrule
The service was not great. & negative \\
Although they tried, the meal was not wonderful. & negative \\
The flight was anything but great. & negative \\
The product was not poor. & positive \\
\bottomrule
\end{tabular}
\caption{Examples from \texttt{negation.csv} (216 rows).}
\end{table}

On this subset DistilBERT (SST-2) fails 34.3\% overall, in the paper's range
of roughly 19\% to 54\%. The aggregate hides where it breaks. Splitting by
pattern, the plain template fails 2.8\% and the buried one 0\%, but
\texttt{anything but} fails 100\%: every single row. The model reads the
positive adjective and ignores the construction completely. One template
isolates the failure, which is the point of testing capabilities one at a
time rather than reporting a single accuracy number.

\subsection{Subset 2 -- Names (INV)}

Here I fix a full sentence and change only a person name or a place name.
None of these names carry sentiment, so the prediction should not move across
a group. Each group shares an id so the evaluation can check that every
variant gets the same label.

\begin{table}[h]
\centering
\begin{tabular}{ll}
\toprule
Group & Text \\
\midrule
0 & Sarah said the service was great. \\
0 & Priya said the service was great. \\
2 & The trip to Berlin was wonderful. \\
2 & The trip to Mumbai was wonderful. \\
\bottomrule
\end{tabular}
\caption{Examples from \texttt{names.csv} (24 rows, 4 groups).}
\end{table}

This is unlabelled on purpose. The test is not whether the prediction is
correct but whether it stays constant when an irrelevant token changes. In
the paper, swapping locations changed predictions for up to about 21\% of
cases, and swapping person names for up to about 15\%.

\subsection{Subset 3 -- Noise (DIR/INV)}

I take a clearly polarised sentence and append a neutral fragment: a handle,
a URL, or a context-free fact. The added text carries no sentiment, so the
prediction should not change direction. Each row keeps both the base and the
perturbed version so the two can be compared directly.

\begin{table}[h]
\centering
\begin{tabular}{ll}
\toprule
Base & Perturbed \\
\midrule
The service was great. & The service was great. @pi9QDK \\
The flight was terrible. & The flight was terrible. https://t.co/PWK1jb \\
The meal was wonderful. & The meal was wonderful. It is Tuesday. \\
\bottomrule
\end{tabular}
\caption{Examples from \texttt{noise.csv} (24 rows, 6 groups).}
\end{table}

The paper runs the same idea by adding random URLs and handles to tweets and
finds that predictions shift for between 7\% and 25\% of inputs. A model that
moves its prediction here is reacting to noise.

\section{How the data is generated}

Everything comes from one short script, \texttt{generate.py}. It holds six
small word lists (positive adjectives, negative adjectives, neutral nouns,
person names, places, filler fragments). The negation subset is the Cartesian
product of nouns and adjectives dropped into each of the three templates. The
names subset
fills four sentence templates with each name or place in turn. The noise
subset pairs each base sentence with each filler fragment. The script writes
three CSV files and prints the row count of each. A second script,
\texttt{evaluate.py}, runs a pretrained sentiment model
(\texttt{distilbert-base-uncased-finetuned-sst-2-english}) over the three
files and prints a failure rate per subset: prediction versus gold for
negation, prediction agreement within a group for names, and base versus
perturbed agreement for noise.

The lists are kept small and hand-checked so that every adjective has an
unambiguous polarity. If the lexicon contained a borderline word, the gold
label for the negation subset would no longer be safe, so curation matters
more than size here.

\section{Links}

\begin{itemize}[noitemsep]
  \item Dataset and code:
    \url{https://github.com/jasraj-jsa/Control-Dataset-for-Behavioural-Testing-of-Sentiment-Models}
  \item Paper: Ribeiro et al.\ (2020),
    \url{https://arxiv.org/abs/2005.04118}
  \item CheckList software:
    \url{https://github.com/marcotcr/checklist}
\end{itemize}

\end{document}
